\documentclass[12pt]{article}
\usepackage[T1]{fontenc}
\usepackage[utf8]{inputenc}
\usepackage{lmodern}
\usepackage{textcomp}
\usepackage{enumitem}
\usepackage{geometry}
\geometry{margin=1in}
\begin{document}
\begin{center}
Answer Key - Geography - K-12 Level Assessment 7 \
\small (Answers are ordered exactly as in the question paper.)
\end{center}

\section*{Multiple Choice}
\begin{enumerate}[label=\arabic*.]
\item \textbf{Answer:} D
\\ \textit{Options:} A) a water divider.; B) a watercourse.; C) an artificial boundary.; D) a natural boundary.
\\ \textit{Note:} The Rio Grande is classified as a natural boundary because it is a physical geographic feature, specifically a river, that serves as a dividing line between the United States and Mexico.
\item \textbf{Answer:} C
\\ \textit{Options:} A) push factors.; B) pull factors.; C) distance decay.; D) migration selectivity.
\\ \textit{Note:} Distance decay refers to the principle that as the distance between two locations increases, the likelihood of interaction or migration between them decreases, making it the correct answer for the tendency of migration to decrease with distance.
\item \textbf{Answer:} D
\\ \textit{Options:} A) zero population growth.; B) rapid growth rate.; C) homeostatic plateau.; D) demographic momentum.
\\ \textit{Note:} Demographic momentum refers to the phenomenon where a population continues to grow even after reaching replacement fertility levels due to the larger proportion of individuals in reproductive age who were born during previous higher birth rates.
\item \textbf{Answer:} D
\\ \textit{Options:} A) It changed centuries-old methods of farming.; B) It was an attempt to feed a growing world population.; C) It centered on increased rice, wheat, and maize yields.; D) It greatly improved crop yields in sub-Saharan Africa.
\\ \textit{Note:} The statement is correct because the Green Revolution primarily benefited countries in Asia and Latin America, but its impact on sub-Saharan Africa was limited due to various challenges such as inadequate infrastructure and access to technology.
\item \textbf{Answer:} D
\\ \textit{Options:} A) They are technologically advanced.; B) They are wealthy.; C) They are powerful.; D) They are less developed.
\\ \textit{Note:} The statement is correct because "core" countries are characterized by advanced economies, high levels of industrialization, and greater political stability, making them more developed compared to "peripheral" or "semi-peripheral" countries.
\end{enumerate}

\section*{True / False}
\begin{enumerate}[label=\arabic*.]
\setcounter{enumi}{5}
\item \textbf{Answer:} False
\\ \textit{Options:} A) True; B) False
\\ \textit{Note:} Urban agriculture often requires substantial water use and can lead to pollution from fertilizers and pesticides, which may negatively impact water quality rather than improve it.
\item \textbf{Answer:} True
\\ \textit{Options:} A) True; B) False
\\ \textit{Note:} The Rhine River serves as a natural boundary between several countries in Europe, including France, Germany, and the Netherlands, due to its geographical location and significant historical role as a border.
\item \textbf{Answer:} True
\\ \textit{Options:} A) True; B) False
\\ \textit{Note:} Bulk-reducing products lose mass during processing, allowing them to be transported longer distances to markets without incurring high transportation costs.
\item \textbf{Answer:} False
\\ \textit{Options:} A) True; B) False
\\ \textit{Note:} Secondary economic activities involve the processing and manufacturing of raw materials rather than the extraction of natural resources, and they are significantly influenced by the physical environment in terms of location, resources, and infrastructure.
\item \textbf{Answer:} False
\\ \textit{Options:} A) True; B) False
\\ \textit{Note:} The statement is false because Dixie is not universally defined as a functional region; rather, it is primarily considered a vernacular region based on cultural and historical identity rather than specific functional interactions or economic ties.
\end{enumerate}

\section*{Long Form Response}
\begin{enumerate}[label=\arabic*.]
\setcounter{enumi}{10}
\item \textbf{Answer:} The Pacific Ring of Fire is a crucial area for global tectonic activity, as it is home to a large number of earthquakes and volcanic eruptions due to the movement of tectonic plates. Countries located along this ring, such as Chile, experience significant geological events because they sit on the boundary of the Nazca and South American plates. This tectonic activity poses risks for Chile, including powerful earthquakes and volcanic eruptions that can impact infrastructure and communities. Additionally, the Ring of Fire affects Chile's landscape, contributing to the creation of mountains and fertile soil, which can be beneficial for agriculture despite the risks involved. Understanding the Pacific Ring of Fire is essential for recognizing how natural forces shape the environment and influence the lives of people living in these areas.
\\ \textit{Note:} correct because it accurately identifies the Pacific Ring of Fire as a region of intense tectonic activity due to plate movements, detailing how countries like Chile are directly impacted by the resulting earthquakes and volcanic eruptions.
\item \textbf{Answer:} Christianity, as a monotheistic religion founded on the teachings of Jesus Christ, has significantly influenced both cultural and geographical landscapes in regions where it is practiced. Culturally, Christianity has shaped art, music, and traditions, leading to the creation of iconic works such as cathedrals and religious paintings, as well as holidays like Christmas and Easter that are celebrated worldwide. Geographically, the spread of Christianity has often been linked to the establishment of churches and schools, which can be seen in the layout of cities and towns, particularly in Europe and the Americas. Additionally, the religion has played a role in the development of social structures and legal systems, reflecting its values in local customs and laws. Overall, Christianity's presence has enriched the cultural identity and influenced the physical landscape of many societies.
\\ \textit{Note:} correct because it identifies the profound cultural influences of Christianity on art, music, and traditions, as well as its geographical spread through the establishment of churches and the observance of holidays, demonstrating its significant impact on both cultural and geographical landscapes.
\end{enumerate}

\vfill
\noindent\textit{End of Answer Key}
\end{document}